% ==========================================================
% == Section 2: Related Work (minimal, 0.5 page)           ==
% == Full Related Work appears in Appendix \ref{apdx:related-work-full}
% ==========================================================
\section{Related Work}
\label{sec:related}

% [段 2.1 · Jailbreaks + static alignment] 大意:极简版攻击演化 + 静态对齐的 prior work,用作 intro 的补充。
% 撰写逻辑:一句话从 DAN / GCG / PAIR / TAP / AutoDAN / AutoDAN-Turbo / X-Teaming / MAJIC 串过;一句话把 RLHF / guardrail / 静态 red-team 数据集作为 static alignment 的典型代表并指出其对预收集分布的依赖。
% 中文对应内容:攻击演化方面,研究路径已从早期人工模板（DAN)经由梯度优化（GCG）、LLM 驱动迭代（PAIR、TAP、AutoDAN、AutoDAN-Turbo)发展到今天的多轮 agentic 攻击（X-Teaming、MAJIC）;对齐方面,RLHF、guardrail 与静态 red-team 数据集（WildGuard、HarmBench、Sorry-Bench）等方法本质上都依赖预收集的攻击分布,在面对分布漂移时结构性失守。详细综述见附录 \ref{apdx:related-work-full}。
\paragraph{LLM jailbreaks and static alignment.} Jailbreak research has progressed from hand-crafted templates~\citep{shen2023doanythingnow} through gradient-based optimization~\citep{zou2023universal}, LLM-driven iterative rewriting~\citep{chao2023pair,mehrotra2024tap,liu2023autodan,liu2024autodan}, to multi-turn and agentic attacks~\citep{rahman2025xteaming,qi2025majic,samvelyan2024rainbow,liu2024flipattack}. Defensive approaches dominated by RLHF~\citep{ouyang2022training,bai2022constitutional,dai2023safe} and fixed-distribution red-teaming data~\citep{jiang2024wildteaming,mazeika2024harmbench,xie2024sorry} are structurally tied to pre-collected attack distributions and thus brittle under drift, motivating the co-evolutionary line below.

% [段 2.2 · Co-Evolutionary Safety Alignment] 大意:简述协同演化先验工作;强调都没系统性解决覆盖 + 记忆两大问题,特别点出 TriPlay-RL 和 SSP 作为最接近的 prior works。
% 中文对应内容:协同演化安全对齐以 Self-RedTeam 的零和自博弈和 MAGIC 的非对称序贯博弈为代表,进一步的变体包括 AdvEvo-MARL、ACE-Safety、SEAS、SSP、TriPlay-RL 以及若干 Stackelberg / 非合作博弈方案。这些工作从均衡概念、多智能体结构等角度推进协同演化,但均未系统解决攻击策略坍缩与防御对抗遗忘;TriPlay-RL 是在协同演化建模中显式引入多样性项的最接近先验工作,但其信号仅为 Self-BLEU + 句向量余弦相似度(表层语义);SSP 维护了回放池并实时更新,但依赖离散 safety score,缺失不确定性建模、偶然性建模与时间衰减。DACE 正面回应这两组结构性不足。
\paragraph{Co-evolutionary safety alignment.} Recent work couples attackers and defenders in a shared RL loop: Self-RedTeam~\citep{liu2025chasing} establishes zero-sum self-play with a Nash equilibrium; MAGIC~\citep{magic2025} casts the interaction as an asymmetric sequential game with a subgame-perfect equilibrium; further variants include AdvEvo-MARL~\citep{pan2025advevo}, ACE-Safety~\citep{acesafety2025}, SEAS~\citep{seas2025}, SSP~\citep{ssp2026}, TriPlay-RL~\citep{triplayrl2026}, and several Stackelberg / non-cooperative formulations~\citep{twoplayergame2024,advgame2025}. None of these address attack strategy collapse and defense adversarial forgetting as joint optimization targets. The two closest prior works occupy complementary positions: TriPlay-RL introduces an explicit diversity term inside a co-evolution loop, but its signal is a Self-BLEU plus sentence-embedding cosine penalty (purely textual/semantic); SSP maintains a replay pool that is refreshed in real time, but updates it from a discrete safety score that cannot represent uncertainty, defender stochasticity, or non-stationarity. \method targets exactly these gaps via a strategy-level coverage reward and a continuous Bayesian replay pool with time decay and Thompson sampling.

% [段 2.3 · Diversity in Red-Teaming] 大意:概述 diversity 两大流派(QD & RL-intrinsic),都在静态目标上;DACE 将之迁移到动态协同。
% 中文对应内容:自动红队领域的多样性研究可归为两类:(i) Quality-Diversity 搜索——Rainbow Teaming 及其加速/扩展变体（Ruby、Ferret、RainbowPlus、QDRT、T-MAP、GRTS、Auto-RT 等),以 MAP-Elites 为核心,针对静态目标构造 archive;(ii) RL 内在多样性奖励——CRT、DiveR-CT、GFlowNet、Jailbreak-R1、CALM、OpenAI RBR 等,通过好奇心/熵/组相对等信号促多样,但仍停留在文本或 embedding 层面。在协同演化中直接移植这些静态或表层机制会遇到两类失败模式（见 §\ref{sec:intro}）;DACE 把多样性重新定位到"策略空间覆盖度",并以归一化 KL 收缩为基础,实现与动态防御者兼容的多样性信号。
\paragraph{Diversity in red-teaming.} Automated red-teaming has explored diversity along two lines: (i) \emph{Quality-Diversity} search---Rainbow Teaming~\citep{samvelyan2024rainbow}, Ruby Teaming~\citep{han2024ruby}, Ferret~\citep{pala2025ferret}, RainbowPlus~\citep{dang2025rainbowplus}, QDRT~\citep{wang2025qdrt}, T-MAP~\citep{guo2025tmap}, GRTS~\citep{grts2023}, Auto-RT~\citep{autort2025}, and the WildTeaming / RedTWIZ pipelines~\citep{jiang2024wildteaming,horal2025redtwiz}---all of which construct behavior-descriptor archives against a \emph{static} target; and (ii) \emph{RL intrinsic-diversity rewards}---CRT~\citep{hong2024curiosity}, DiveR-CT~\citep{zhao2025diverct}, GFlowNet-based attackers~\citep{lee2025gflownet}, Jailbreak-R1~\citep{wang2025jailbreakr1}, CALM~\citep{zheng2025calm}, RedTopic~\citep{redtopic2025}, OpenAI RBR~\citep{beutel2024openaidiverse}, and Active Attacks~\citep{activeattacks2025}---which operate at the token or embedding level. \method departs from both by placing diversity at the \emph{strategy level} inside a dynamic co-evolution loop, and by deriving its reward signal from normalized KL contraction rather than from either static archives or textual novelty. An expanded discussion, including continual-learning analogies for our replay mechanism, is provided in Appendix~\ref{apdx:related-work-full}.
